\title{xLSTM: Extended Long Short-Term Memory}

\begin{abstract}
The introduction of the constant error carousel and gating mechanisms in the 1990s laid the groundwork for the development of the Long Short-Term Memory (LSTM) architecture. Over subsequent decades, LSTMs have remained fundamental components of deep learning, particularly as the core of initial Large Language Models (LLMs). Nonetheless, the emergence of Transformer architectures, distinguished by their efficient self-attention mechanisms conducive to parallelism, has since redefined the benchmarks for large-scale modeling, surpassing the performance of LSTMs. In this work, we pose a straightforward inquiry: What advancements are possible in language modeling when LSTMs are scaled to billions of parameters by utilizing contemporary LLM methodologies, while strategically addressing known deficiencies inherent to LSTMs? Our contributions are twofold. First, we propose exponential gating equipped with tailored normalization and stabilization strategies. Second, we redesign the LSTM’s memory dynamics, resulting in two novel variants: (i) sLSTM, characterized by a scalar memory and update along with an enhanced scheme for memory combination, and (ii) mLSTM, which replaces vector memory with matrix memory and incorporates a covariance-based update, facilitating complete parallelization. By embedding these advanced LSTM structures within residual block frameworks, we construct xLSTM blocks, which can be composited in a residual manner to form full xLSTM architectures. Empirical results indicate that the integration of exponential gating and restructured memory allows xLSTM models to achieve competitive—or even superior—performance and scaling properties relative to prevailing Transformer-based and State Space Models.\\
Code available at: \url{https://github.com/NX-AI/xlstm}
\end{abstract}

\section{Introduction}
\label{sec:intro}

The concepts behind Long Short-Term Memory (LSTM)~\citep{Hochreiter:91,Hochreiter:97nips,Hochreiter:97}, specifically the constant error carousel and the use of gating mechanisms, were proposed as a solution to address the vanishing gradient issue found in recurrent neural networks~\citep{Hochreiter:91,Hochreiter:00book}:

\begin{align}
\label{eq:lstm_idea}
  c_t \ = \ f_t \ c_{t-1} \ + \ 
          i_t \  z_t \ , \quad  
          h_t \ = \ o_t \ \psi({c_t}) \ . 
\end{align}

The constant error carousel involves updating the cell state $c_{t-1}$ (green) additively with the cell inputs $z_t$, with this process being regulated by the sigmoid gates (blue). The input gate $i_t$ along with the forget gate $f_t$ determines how this update unfolds, whereas the output gate $o_t$ is responsible for generating the memory cell's output, that is, the hidden state $h_t$. Before producing the hidden state, the cell state undergoes normalization or squashing by $\psi$, followed by modulation by the output gate.

LSTMs have achieved notable success across a broad range of fields \citep{Hochreiter:01,Hochreiter:07,Schmidhuber:15}, dominating text generation tasks prior to the emergence of Transformers in 2017~\citep{Vaswani:17}. Their capabilities have been well-validated in numerous sequential tasks, including but not limited to text synthesis~\citep{Graves:13,Karpathy:15blog}, handwriting generation~\citep{Graves:13}, sequence-to-sequence translation~\citep{Sutskever:14nips}, program evaluation~\citep{Zaremba:14arxiv}, automatic image captioning~\citep{Karpathy:15,Hossain:19}, code generation~\citep{Karpathy:15blog}, modeling rainfall-runoff dynamics~\citep{Kratzert:18,Kratzert:19}, and hydrological flood forecasting~\citep{Nearing:24}. In the context of reinforcement learning, LSTMs have demonstrated state-of-the-art performance, as exemplified by their integration in systems such as the AlphaStar agent for StarCraft II~\citep{Vinyals:17short}, OpenAI Five for Dota 2~\citep{OpenAI:19}, and magnetic control systems for nuclear fusion~\citep{Degrave:22short}. LSTMs are particularly strong in forming higher-level abstractions—proficiently capturing semantic content within their memory components~\citep{Karpathy:15blog}—a trait illustrated by the presence of specialized number and syntax neurons~\citep{Lakretz:19}, linguistic neurons~\citep{Bau:19}, and sentiment neurons~\citep{Radford:17arxiv}. To this day, LSTMs continue to play a vital role in cutting-edge applications~\citep{Degrave:22short,Nearing:24}, attesting to their enduring significance.

Although LSTMs have achieved notable results, they exhibit three principal drawbacks: (i) Inability to revise stored information. This is illustrated with the {\it Nearest Neighbor Search} problem (refer also to Appendix~\ref{sec:appNearestNeighborSearch}): Given a reference vector, the task requires sequentially traversing a sequence to find the closest vector and output its associated value at the end. The left panel of Figure~\ref{fig:lstmProblems} reports the mean squared error on this task. LSTMs have difficulty updating a previously stored value upon encountering a more similar vector, whereas the proposed xLSTM, with its exponential gating mechanism, effectively addresses this shortfall. (ii) Restricted storage capacity, as the information must be compressed into single-dimensional cell states. This is demonstrated using the {\it Rare Token Prediction} task. The right panel of Figure~\ref{fig:lstmProblems} displays token prediction perplexity from the Wikitext-103 dataset~\citep{Merity:17}, segmented by token frequency. LSTMs show degraded performance on infrequently seen tokens due to their constrained memory, while our xLSTM, leveraging a matrix-based memory, overcomes this issue. (iii) Limited parallelizability stemming from memory interdependencies, specifically the hidden-to-hidden state transitions across time steps, which enforce a strictly sequential computation process.

The shortcomings of LSTM architectures have led to the rise of Transformers~\citep{Vaswani:17} within the field of language modeling. An intriguing question arises: by addressing these constraints and expanding LSTMs to match the scale of modern Large Language Models, what levels of performance become attainable in language modeling tasks?

\section{Extended Long Short-Term Memory}
\label{sec:xLSTM}

In response to the recognized constraints of LSTM, Extended Long Short-Term Memory (xLSTM) incorporates two primary innovations building on the LSTM framework as in Equation~\eqref{eq:lstm_idea}. These enhancements --- namely, exponential gating and distinct new memory mechanisms --- augment the LSTM class by introducing (i) sLSTM (refer to Section~\ref{sec:sLSTM}), characterized by a scalar memory, scalar updates, and memory mixing; and (ii) mLSTM (see Section~\ref{sec:mLSTM}), which utilizes a matrix memory together with a covariance-based (outer product) update strategy that allows for full parallelization. Exponential gating strengthens both sLSTM and mLSTM. The mLSTM achieves parallel computation by omitting memory mixing, or the hidden-to-hidden recurrent interactions. Both variants, mLSTM and sLSTM, are further extensible to multiple memory cells, where the sLSTM allows for mixing among these cells. Additionally, in the sLSTM, it is possible to employ multiple heads in which memory mixing is restricted to within each head’s cells, excluding mixing between different heads, which gives rise to a novel form of memory interaction when combined with exponential gating. In contrast, for mLSTM, there is an equivalence between the roles of multiple heads and multiple memory cells.

The incorporation of these novel LSTM variants within residual block modules yields xLSTM blocks (refer to Section~\ref{sec:xLSTMblock}). When these xLSTM blocks are stacked using residual connections, they form xLSTM architectures (refer to Section~\ref{sec:xLSTMarchitecure}). Figure~\ref{fig:xlstm_sketch} illustrates the structure of the xLSTM architecture along with its constituent parts.

\subsection{Review of the Long Short-Term Memory}
\label{sec:orgLSTM}

The initial concept of LSTM~\citep{Hochreiter:91,Hochreiter:97nips,Hochreiter:97} proposed a scalar memory cell as the core unit for both computation and memory, designed specifically to circumvent the issue of vanishing gradients~\citep{Hochreiter:91,Hochreiter:00book} using a mechanism known as the constant error carousel, which relies on cell state updates. The architecture of the memory cell includes three distinct gates: the input gate, the output gate, and the forget gate—the latter of which was subsequently added by \citet{Gers:00}. The dynamics governing the evolution of the LSTM memory cell at any given time step $t$ can be described as follows:

\begin{align}
c_t \ &= \  f_t \ c_{t-1} \ + \ i_t \ z_t &  & &\text{cell state} \\
h_t  \ &= \ o_t \ \tilde{h}_t \ , & \tilde{h}_t \ &= \ \psi \left( c_t\right)
  &\text{hidden state} \\
z_t \ &= \ \varphi \left( \tilde{z}_t \right) \ , 
  &\tilde{z}_t \ &=  \ \mathbf{w}^\top_{z} \ \mathbf{x}_t \ + \
  r_{z}  h_{t-1} \ + \  b_{z} \ \
  &\text{cell input} \\
i_t \ &= \ \sigma \left( \tilde{i}_t  \right) \ , 
  &\tilde{i}_t \ &= \ \mathbf{w}^\top_{i} \ \mathbf{x}_t \ + \
  r_{i}  \ h_{t-1} \ + \  b_{i} \ \
  &\text{input gate} \\
f_t \ &= \ \sigma \left(  \tilde{f}_t \right) \ , 
  &\tilde{f}_t \ &= \ \mathbf{w}^\top_{f} \ \mathbf{x}_t  \ + \
  r_{f}  \ h_{t-1} \ + \  b_{f} \ \
  &\text{forget gate} \\
o_t \ &= \ \sigma \left( \tilde{o}_t \right) \ , 
  &\tilde{o}_t  \ &= \ \mathbf{w}^\top_{o} \ \mathbf{x}_t \ + \
  r_{o}  \ h_{t-1} \ + \  b_{o} \ \
  &\text{output gate} 
\end{align}

\textit{(Note: The LaTeX equations provided already concisely represent mathematical relations. The paraphrasing only applies to explanatory text; thus, equations and accompanying labels have been kept unchanged as per instructions.)}

The weight vectors $\mathbf{w}_{z}$, $\mathbf{w}_{i}$, $\mathbf{w}_{f}$, and $\mathbf{w}_{o}$ are assigned to the connections from the input vector $\mathbf{x}_t$ to the cell input, the input gate, the forget gate, and the output gate, respectively. The parameters $r_{z}$, $r_{i}$, $r_{f}$, and $r_{o}$ denote recurrent weights linking the previous hidden state $h_{t-1}$ to the cell input and the respective gates. The associated bias terms for each pathway are represented as $b_{z}$, $b_{i}$, $b_{f}$, and $b_{o}$. The functions $\varphi$ and $\psi$ govern activation at the cell input and hidden state (commonly chosen as $\tanh$); $\psi$ particularly serves to bound the values of the cell state, which otherwise could grow without limit. Each gate employs a sigmoid activation, specified by $\sigma \left( x \right)= 1/(1+\exp(-x))$. In subsequent variants, vectors of scalar memory cells $\mathbf{c}_t \in \mathbb{R}^d$ were introduced, facilitating the implementation of recurrent weight matrices $\mathbf{R} \in \mathbb{R}^{d \times d}$ to mix outputs across memory cells~\citep{Greff:15} (refer to Appendix~\ref{sec:apporgLSTM} for further discussion). Ablation analyses revealed that the integrity of every memory cell component is essential~\citep{Greff:15}.

\subsection{sLSTM}
\label{sec:sLSTM}

To enable LSTMs to modify their storage decisions, we incorporate exponential gates (illustrated in red) along with mechanisms for normalization and stabilization. Specifically, both input and forget gates are permitted to utilize exponential activation functions. For the purpose of normalization, we propose a normalizer state which accumulates the sum of the input gate multiplied by each subsequent forget gate.

The scalar sLSTM forward pass is:

\begin{align}
\label{eq:slstmforward}
c_t \ &= \  f_t \ c_{t-1} \ + \ i_t \ z_t & & 
 &\text{cell state} \\
n_t \ &= \  f_t \ n_{t-1} \ + \ i_t  & & 
 &\text{normalizer state} \\
h_t \ &= \ o_t \ \tilde{h}_t \ ,
  &\tilde{h}_t \ &= \ c_t / n_t
&\text{hidden state} \\
z_t \ &= \ \varphi \left( \tilde{z}_t \right) \ , 
  &\tilde{z}_t \ &=  \ \mathbf{w}^\top_{z} \ \mathbf{x}_t \ + \
  r_{z}  h_{t-1} \ + \  b_{z}
  &\text{cell input} \\
i_t \ &= \ \!\exp \left( \tilde{i}_t  \right) \ , 
  &\tilde{i}_t \ &= \ \mathbf{w}^\top_{i} \ \mathbf{x}_t \ + \
  r_{i}  \ h_{t-1} \ + \  b_{i}
  &\text{input gate} \\
f_t \ &= \ \sigma \left(  \tilde{f}_t \right) \ \text{OR} \
\!\exp \left(  \tilde{f}_t \right) \ ,
  &\tilde{f}_t \ &= \ \mathbf{w}^\top_{f} \ \mathbf{x}_t  \ + \
  r_{f}  \ h_{t-1} \ + \  b_{f}
  &\text{forget gate} \\
o_t \ &= \ \sigma \left( \tilde{o}_t \right) \ , 
  &\tilde{o}_t  \ &= \ \mathbf{w}^\top_{o} \ \mathbf{x}_t \ + \
  r_{o}  \ h_{t-1} \ + \  b_{o}
  &\text{output gate} 
\end{align}

We adapt the classic LSTM gating mechanisms—specifically, those relying on dependencies on inputs and/or hidden states along with a bias component—to these novel architectural designs. As exponential activation functions are susceptible to producing exceedingly large outputs that risk numerical overflow, we incorporate an auxiliary state \mbox{$m_t$~\citep{Milakov:18arxiv}} to provide additional stability to the gates:

\begin{align}
\label{eq:slstmstabil}
m_t \ &= \ \max \left( \log ( f_t ) + m_{t-1} , \log ( i_t ) \right) &\text{stabilizer state} \\
i'_t \ &= \ \!\exp \left( \log \left ( i_t \right) - m_t \right) \ = \ \!\exp \left( \tilde{i}_t - m_t \right) \ \, 
  &\text{stabil. input gate} \\
f'_t \ &= \ \!\exp \left( \log \left( f_t \right) + m_{t-1} - m_t \right) \ \, 
  &\text{stabil. forget gate}
\end{align}

As demonstrated in Appendix~\ref{sec:appsLSTM}, substituting $f_t$ with $f'_t$ and $i_t$ with $i'_t$ during the forward computation does not alter either the overall network output or the gradients of the loss function with respect to the parameters.

\paragraph{New Memory Mixing.} Similar to the original LSTM, sLSTM can incorporate several memory cells (refer to Appendix~\ref{sec:appsLSTM}).
Having multiple memory cells allows for memory mixing through recurrent pathways $\mathbf{R}_{\mathbf{z}}$, $\mathbf{R}_{\mathbf{i}}$, $\mathbf{R}_{\mathbf{f}}$, $\mathbf{R}_{\mathbf{o}}$, which connect the hidden state vector $\mathbf{h}$ to the memory cell input $\mathbf{z}$ as well as to the gating units $\mathbf{i}$, $\mathbf{f}$, and $\mathbf{o}$, respectively. A distinctive feature in this mixing process is introduced by exponential gating. In the updated sLSTM, each head can possess several memory cells with internal memory mixing, although mixing does not occur between different heads. By integrating heads within the sLSTM and employing exponential gating, this model brings forth a novel framework for combining memory states.

\subsection{mLSTM}
\label{sec:mLSTM}

In order to improve the storage capabilities of LSTMs, we generalize the memory cell from a single scalar value $c \in \mathbb{R}$ to a full matrix representation $\mathbf{C} \in \mathbb{R}^{d \times d}$. With this modification, memory retrieval relies on matrix multiplication. Specifically, at each timestep $t$, the model stores a vector pair: a key $\mathbf{k}_t \in \mathbb{R}^d$ and a value $\mathbf{v}_t \in \mathbb{R}^d$, adopting terminology from the Transformer framework. Subsequently, at a future time $t+\tau$, the value $\mathbf{v}_t$ can be retrieved using a query vector $\mathbf{q}_{t+\tau} \in \mathbb{R}^d$. This approach corresponds to the structure of Bidirectional Associative Memories (BAMs) as described in \citep{Kohonen:72,Anderson:72,Nakano:72,Anderson:77}.

The covariance update rule~\citep{Sejnowski:77,Dayan:91} for storing a key-value pair is

\begin{align}
\mathbf{C}_t \ &= \ \mathbf{C}_{t-1} \ + \ \mathbf{v}_{t} \ \mathbf{k}_t^\top \ .
\end{align}

We consider a layer-norm operation applied prior to the input projection for keys and values, ensuring that their means are centered at zero. According to~\citep{Dayan:91}, the covariance update procedure yields an optimal strategy for maximizing the separability of retrieved binary vectors, which also maximizes the signal-to-noise ratio. Enhanced separability is achievable if retrieval is restricted to pairwise interactions, although this incurs quadratic computational complexity, as is the case in attention mechanisms~\citep{Krotov:16,Krotov:17,Ramsauer:21}. Notably, the covariance update coincides with the approach used by Fast Weight Programmers \citep{Schmidhuber:92ncfastweights,Schlag:21}. This framework has been subsequently modified by incorporating a constant decay factor applied to $\mathbf{C}_{t-1}$ and a constant learning rate for the term 
$\mathbf{v}_{t} \mathbf{k}_t ^\top$~\citep{Ba:16}. Drawing on these developments, we embed the covariance update process within the LSTM architecture, identifying the forget gate as the mechanism for applying decay, the input gate as modulating the learning rate, and the output gate as a scale for the retrieved output.

In the case of this matrix memory, the normalizer state is defined as a weighted combination of the key vectors, with each key vector being scaled by the corresponding input gate along with the product of all subsequent forget gates. This ensures that the normalizer state reflects the cumulative effect of the gating mechanisms. Because the dot product between the query and the normalizer state might approach zero, we apply the absolute value to this dot product and ensure it is at least as large as a predetermined threshold (typically 1.0), consistent with previous work~\citep{Sun:23arxiv}. The forward computation for mLSTM is as follows:

\begin{align}
\label{eq:mlstm_recurrent_begin}
\mathbf{C}_t \ &= \  f_t \ \mathbf{C}_{t-1} \ + \ 
  i_t \ \mathbf{v}_t \ \mathbf{k}_t^\top & &  &\text{cell state} \\
\mathbf{n}_t \ &= \  f_t \ \mathbf{n}_{t-1} \ + \ 
  i_t \ \mathbf{k}_t & &  &\text{normalizer state} \\
\mathbf{h}_t  \ &= \ \mathbf{o}_t \ \odot \ \tilde{\mathbf{h}}_t
\ , \qquad \qquad 
\tilde{\mathbf{h}}_t \ = \   \mathbf{C}_t \mathbf{q}_t \ / \ 
  \max \left\{ |\mathbf{n}_t^\top \mathbf{q}_t|, 1 \right\} 
   &&&\text{hidden state} \\
\mathbf{q}_t \ &= \ \mathbf{W}_q \ \mathbf{x}_t \ + \ \mathbf{b}_q  & &  &\text{query input} \\
\mathbf{k}_t \ &= \ \frac{1}{\sqrt{d}} \mathbf{W}_k \ \mathbf{x}_t \ + \ \mathbf{b}_k  & &  &\text{key input} \\
\mathbf{v}_t \ &= \ \mathbf{W}_v \ \mathbf{x}_t \ + \ \mathbf{b}_v  & &  &\text{value input} \\
i_t \ &= \ \!\exp \!  \! \left( \tilde{i}_t  \right) \ , 
 \qquad \qquad \ \ \ \ \,
  \tilde{i}_t \ = \ \mathbf{w}^\top_{i} \ \mathbf{x}_t \ + \  b_{i} \ \ \ 
  &&&\text{input gate} \\
f_t \ &= \ \sigma \!  \left(  \tilde{f}_t \right) \ \text{OR} \ \!\exp \!  \! \left(  \tilde{f}_t \right) , \ \ \: \!
\tilde{f}_t \ = \ \mathbf{w}^\top_{f} \ \mathbf{x}_t  \ + \
  b_{f}
  &&& \text{forget gate} \\
\label{eq:mlstm_recurrent_end}
\mathbf{o}_t \ &= \ \sigma \left( \tilde{\mathbf{o}}_t \right) \ , \qquad \qquad \quad \ \ \ \,
  \tilde{\mathbf{o}}_t  \ = \ \mathbf{W}_{\mathbf{o}} \ \mathbf{x}_t \ + \
  \mathbf{b}_{\mathbf{o}}
  &&&\text{output gate} 
\end{align}

The mLSTM architecture, similar to the standard LSTM, is capable of supporting multiple memory cells. In mLSTM, there is no memory mixing mechanism, making the use of multiple heads functionally equivalent to employing multiple cells. To ensure stability in the exponential gating mechanism within mLSTM, we adopt the same stabilization strategies as those implemented for sLSTM (refer to Equation~\ref{eq:slstmstabil}). Given the absence of memory mixing in mLSTM, its recurrence relations can be reformulated for parallel execution. Additional information can be found in Appendix~\ref{sec:appmLSTM}.

\subsection{xLSTM Architecture}
\label{sec:xLSTMarch}

\paragraph{xLSTM Blocks.}
\label{sec:xLSTMblock}
The function of an xLSTM block is to generate a non-linear encoding of the sequence's historical data in a high-dimensional representation, thereby facilitating a clearer discrimination between distinct contexts or past trajectories. Achieving such a separation is essential for accurately forecasting the subsequent element in a sequence, such as the next token. Our approach draws upon Cover’s Theorem~\citep{Cover:65}, which posits that non-linear mappings to higher-dimensional spaces tend to make patterns more easily separable by linear functions compared to their separability in the original space. We analyze two types of residual block designs: (i) A residual block featuring post up-projection (akin to the Transformer architecture), in which the history is non-linearly encoded within the original space, subsequently projected via a linear transformation into a higher-dimensional space, subjected to a non-linear activation, and then projected back linearly to the initial dimensionality; this configuration is illustrated in the left panel of Figure~\ref{fig:Backbones} and in the third column of Figure~\ref{fig:xlstm_sketch}. An expanded depiction is also available in Figure~\ref{fig:appsLSTM_detailed} within the appendix. (ii) Alternatively, a residual block may utilize a pre up-projection strategy (similar to State Space Models), where input is first projected linearly into a higher-dimensional space,

The model performs a nonlinear condensation of historical information within a high-dimensional feature space and subsequently applies a linear mapping to return this summary to the original representation space. For xLSTM blocks that include an sLSTM component, we implement a post up-projection mechanism. In contrast, for xLSTM blocks with an mLSTM module, we employ a pre up-projection step, as the expanded memory capacity operates within the higher-dimensional context. More comprehensive explanations can be found by consulting the left side of Figure~\ref{fig:Backbones}, the third column of Figure~\ref{fig:xlstm_sketch}, or Figure~\ref{fig:appsLSTM_detailed} in the appendix.

\paragraph{xLSTM Architecture. }
\label{sec:xLSTMarchitecure}
The xLSTM architecture is developed through the residual stacking of modular components~\citep{Srivastava:15,He:16}. We adopt the widely utilized pre-LayerNorm~\citep{Ba:16b} residual frameworks, which serve as the backbone for modern Large Language Models. An illustration of this structure is provided in the final column of Figure~\ref{fig:xlstm_sketch}.

\subsection{Memory and Speed Considerations}

In contrast to Transformers, xLSTM networks feature linear computational requirements and maintain fixed memory complexity as sequence length increases. Owing to the compressive nature of xLSTM memory, these networks are particularly advantageous for use in industrial contexts and edge-device deployments.

mLSTM’s memory operates without additional parameters; however, its $d \times d$ matrix storage and corresponding update procedures are computationally intensive. Our approach involves balancing the trade-off between increasing memory capacity and the resulting computational demands. Nonetheless, these operations are amenable to parallel processing on GPUs, thus contributing minimally to overall elapsed time during execution.

Although mLSTM, similar to FlashAttention~\citep{Dao:22,Dao:23} and GLA~\citep{Yang:23arxiv}, supports parallel execution, sLSTM lacks this capability owing to its memory mixing through hidden-to-hidden interactions. Nonetheless, we introduced an efficient CUDA-based solution with GPU memory optimizations down to the register level, resulting in a runtime overhead of generally less than a factor of two compared to mLSTM.

\section{Related Work}
\label{sec:related}

\paragraph{Linear Attention.} Various approaches have been developed to address the quadratic computational overhead of Transformer attention in relation to context length, enabling a linear scaling with input size. The Synthesizer architecture produces synthetic attention weights independent of direct token-to-token comparisons~\citep{Tay:20arxiv}. Linformer leverages low-rank projections to approximate the self-attention mechanism efficiently~\citep{Wang:20arxiv}. Linear Transformer modifies the standard attention formulation to achieve linear time complexity~\citep{Katharopoulos:20}. Performer employs a positive orthogonal random features method to achieve a linear approximation of the attention softmax operation~\citep{Choromanski:21}. In addition, attention has been substituted by high-speed, long-range convolutional layers in Structured Global Convolution (SGConv)~\citep{Li:22} and in the Hyena Hierarchy architecture~\citep{Poli:23}.

\paragraph{State Space Models.} State Space Models (SSMs) have gained significant attention in recent times due to their linear scaling with input length and competitive performance when benchmarked against Transformers. The emergence of the Structured State Space sequence model (S4)~\citep{Gu:21} marked an initial breakthrough, which was subsequently expanded upon by the Diagonal State Space (DSS) model~\citep{Gupta:22}, the Gated State Space (GSS) models~\citep{Mehta:22}, S5 model~\citep{Smith:22}, Bidirectional Gated SSM (BiGS)~\citep{Wang:22}, H3 model~\citep{Fu:23}, and most recently, Mamba~\citep{Gu:24arxiv}.

\paragraph{Recurrent Neural Networks.} Recurrent Neural Networks (RNNs) have been proposed as alternatives to Transformer architectures and attention mechanisms, primarily due to their computational linearity with respect to sequence length. RNNs featuring Deep Linear Recurrent Units (LRUs) have achieved strong performance in language modeling tasks~\citep{Orvieto:23, De:24arxiv}. Similarly, the Hierarchically Gated Linear RNN (HGRN)~\citep{Qin:23} and its successor, HGRN2~\citep{Qin:24arxiv}, have demonstrated notable results. Another prominent RNN-based method for large language models, RWKV~\citep{Peng:23arxivshort,Peng:24arxivshort}, has delivered results comparable to those of Transformer models.

\paragraph{Gating.}
Gating represents a foundational concept introduced in LSTM, and has subsequently been independently rediscovered and re-envisioned within numerous contemporary methodologies. A variety of models have integrated gating into their architectures, including HGRN~\citep{Qin:23}, HGRN2~\citep{Qin:24arxiv}, Gated Linear Attention (GLA)~\citep{Yang:23arxiv}, Gated State Space (GSS) models~\citep{Mehta:22}, Bidirectional Gated SSM (BiGS)~\citep{Wang:22}, Moving Average Equipped Gated Attention (MEGA)~\citep{Ma:22}, RWKV~\citep{Peng:23arxivshort}, as well as Mamba~\citep{Gu:24arxiv}.

\paragraph{Covariance Update Rule.} In order to boost storage capabilities, we introduced a matrix memory to the mLSTM cell that leverages a covariance update rule. Various approaches have adopted similar update strategies, such as Fast Weight Programmers \citep{Schmidhuber:92ncfastweights,Schlag:21}, RWKV-5 and RWKV-6~\citep{Peng:24arxivshort}, Retention~\citep{Sun:23arxiv}, Linear Transformer~\citep{Katharopoulos:20}, and HGRN2~\citep{Qin:24arxiv}.

\paragraph{Most Related.} In terms of conceptual similarities, xLSTM is most akin to models such as Retention~\citep{Sun:23arxiv}, RWKV~\citep{Peng:23arxivshort,Peng:24arxivshort}, and HGRN2~\citep{Qin:24arxiv}, all of which utilize the ideas of matrix memory and gating mechanisms. Despite these shared characteristics, unlike the recently proposed sLSTM, these earlier models lack support for memory mixing. The introduction of memory mixing is particularly important for addressing state tracking challenges, which makes LSTMs inherently more powerful and expressive than both State Space Models (SSMs) and Transformers~\citep{Merrill:24,Deletang:23}. Accurate state tracking is critical for tasks such as code evaluation or maintaining awareness of entities throughout lengthy narratives.

\paragraph{Residually Stacking Architectures.} In alignment with the design principles underpinning nearly all modern large-scale deep learning models, xLSTM architectures employ residual stacking of modular components~\citep{Srivastava:15,He:16}.

This strategy was foundational in the development of deep convolutional networks~\citep{He:16} and has also become central to the architecture of Transformers~\citep{Vaswani:17}. The latter, in particular, underpin the capabilities of Large Language Models (LLMs) such as GPT-3~\citep{Brown:20short}, ChatGPT~\citep{Schulman:22short}, GPT-4~\citep{Achiam:23gpt4short}, Megatron-LM~\citep{Shoeybi:19}, Gopher~\citep{Rae:21short}, ERNIE 3.0 Titan~\citep{Wang:21arxivshort}, GLaM ~\citep{Du:21glamshort}, Chinese M6~\citep{Lin:21}, the multilingual AlexaTM 20B~\citep{Soltan:22}, OPT~\citep{Zhang:22opt}, Chinchilla~\citep{Hoffmann:22short}, BLOOM~\citep{Scao:22arxivshort}, GLM-130B~\citep{Zeng:22arxivshort}, LaMDA~\citep{Thoppilan:22short}, PaLM~\citep{Chowdhery:22short}, Llama~\citep{Touvron:23arxiv}, and Gemini~\citep{GeminiTeam:23arxiv,Reid:24arxiv}.

\section{Experiments}
\label{sec:Exp}

This section presents a thorough empirical assessment of xLSTM in the context of language modeling, benchmarking it against established approaches. Section~\ref{sec:ExpSyntethic} explores xLSTM’s unique strengths through synthetic benchmarks. In Section~\ref{sec:ExpComparison}, we benchmark multiple contemporary language modeling architectures, evaluating their validation set perplexity after training on a 15B token sample from SlimPajama~\citep{Soboleva:23}. Ablation experiments targeting xLSTM are also conducted using the same data source. Furthermore, we investigate the scaling patterns of each approach, following the methodology outlined by \citet{Kaplan:20} and \citet{Brown:20short}.

Section~\ref{sec:ExpLanguage} details a comprehensive language modeling study, training xLSTM and the top competitors identified in Section~\ref{sec:ExpComparison} on a substantially larger dataset consisting of 300B tokens from SlimPajama~\citep{Soboleva:23}. The analysis encompasses several aspects: evaluating extrapolation capacity to extended contexts, validation set perplexity and downstream task performance as described by~\citep{Sutawika:24}, performance across 571 specialized text domains provided by the PALOMA language benchmark~\citep{Magnusson:23arxivshort}, and an updated examination of scaling properties in the large-data regime—with the dataset increased by a factor of 20.

Throughout all experiments, we denote xLSTM[$a$:$b$] as the configuration where the proportion of mLSTM-based to sLSTM-based xLSTM blocks is $a/b$. As an illustration, xLSTM[7:1] signifies a setup in which seven out of eight blocks are based on mLSTM, while one block uses sLSTM. With the standard total of 48 blocks, this corresponds to 42 mLSTM-based and 6 sLSTM-based blocks. Additionally, every experimental setting includes the application of up-projection blocks both before and after, for mLSTM and sLSTM, respectively.

\subsection{Synthetic Tasks and Long Range Arena}
\label{sec:ExpSyntethic}
To begin, we evaluate how well the xLSTM's novel exponential gating mechanism with memory mixing performs on formal languages~\citep{Deletang:23}. Next, we examine the capabilities of xLSTM's matrix memory on the Multi-Query Associative Recall benchmark~\citep{Arora:23arxiv}. Lastly, we investigate the model’s proficiency with long sequence processing by benchmarking it on Long Range Arena tasks~\citep{Tay:21}.

\paragraph{Evaluation of xLSTM’s Exponential Gating Mechanism with Memory Mixing.} We assess the efficacy of xLSTM’s innovative exponential gating mechanism integrated with memory mixing, which is intended to allow effective state tracking in sequential tasks~\citep{Merrill:24, Merrill:23}. We build upon and expand the formal language benchmarks introduced by \citet{Deletang:23} to accommodate multi-length training protocols, facilitating length extrapolation studies. For comprehensive information regarding the tasks and additional experimental outcomes, refer to Appendix~\ref{sec:appExpSynthetic-formalLang}. In our study, we benchmark xLSTM against a variety of architectures such as Transformers, State Space Models, and Recurrent Neural Networks. Model performance is assessed based only on those tokens deemed crucial for the given task, with scores normalized to range from 0 (chance-level performance) to 1 (perfect accuracy). Specifically, we conduct comparisons among two-layer configurations for each of the following: xLSTM[0:1] (which is functionally equivalent to sLSTM), xLSTM[1:0] (equivalent to mLSTM), xLSTM[1:1], Llama, Mamba, RWKV, Retention, Hyena, LSTM, and LSTM components within Transformer blocks (LSTM (Block)). A summary of the findings is depicted in Figure~\ref{fig:formal-main}. Notably, architectures such as Transformers and State Space Models lacking memory mixing, and therefore unable to track state, are incapable of addressing even simple regular languages, such as the parity problem. This aligns with previously reported results indicating the inherent limitations of Transformers and State Space Models in comparison with Recurrent Neural Networks~\citep{Merrill:24, Merrill:23, Deletang:23}.

\paragraph{Test of xLSTM's Memory Capacities on Associative Recall Tasks.} This study evaluates the matrix memory mechanism of xLSTM by measuring its ability to store and retrieve information on the Multi-Query Associative Recall task~\citep{Arora:23arxiv}: In this task, sequences are constructed using key-value pairs selected at random from an extensive vocabulary; the model is required to memorize these pairs for subsequent retrieval. To increase the complexity beyond the original formulation, we scale the task up to 256 key-value pairs and extend the context window as far as 2048. This permits a more comprehensive assessment of the memory capabilities across various architectures. We benchmark 2-block model variants of Llama, Mamba, RWKV-5, RWKV-6, xLSTM[1:1], and xLSTM[1:0]. Model performance is quantified by the accuracy with which key-value pairs are recalled. Given that Transformer-based models (such as Llama) exhibit memory scaling exponentially with respect to the representation dimension~\citep{Ramsauer:21}, they are considered the optimal baseline for this benchmark. Figure~\ref{fig:mqar-main} presents the findings.

Among non-Transformer alternatives, xLSTM[1:1] demonstrates superior recall performance, holding this advantage even for smaller-scale models. Notably, the inclusion of the sLSTM block does not impair, but rather augments, the model's memory capacity—a benefit particularly pronounced under the highest challenge level of 256 key-value pairs. Further outcomes, including analyses on model behavior beyond the training context lengths, are discussed in Appendix~\ref{sec:appExpSynthetic-mqar}, where it is shown that xLSTM’s improved memory architecture supports robust extrapolation.

\paragraph{Test of xLSTM's Long Context Capabilities on Long Range Arena.} To evaluate how xLSTM manages extended sequences and substantial contextual information, we benchmark it against other approaches using the Long Range Arena~\citep{Tay:21}. Across the full suite of tasks, xLSTM reliably achieves high results, indicating its architecture is particularly well-suited for addressing the demands of long context scenarios.

For more details, see Appendix~\ref{sec:appExpSynthetic-lra}.

\subsection{Method Comparison and Ablation Study}
\label{sec:ExpComparison}

To explore our central research question—namely, the capabilities of our proposed LSTM variants in large-scale language modeling—we conduct experiments by training xLSTMs, Transformers, State Space Models, and additional approaches using 15B tokens from the SlimPajama dataset, all within an auto-regressive framework. The resulting models are evaluated on a validation set, and we also undertake ablation studies to analyze the components of xLSTMs.

\paragraph{Comparative Analysis of xLSTM and Baseline Approaches.} To facilitate a systematic comparison, we train models on a corpus of 15B tokens derived from SlimPajama~\citep{Soboleva:23}, and subsequently assess their performance via perplexity measured on a designated validation set. The suite of approaches under evaluation encompasses: xLSTM (the newly proposed model), GPT-3 (Transformer architecture)~\citep{Brown:20short}, Llama (Transformer)~\citep{Touvron:23arxiv}, H3 (an SSM-based method)~\citep{Fu:23}, Mamba (SSM)~\citep{Gu:24arxiv}, RWKV-4, RWKV-5, RWKV-6 (all RNN-based models)~\citep{Peng:23arxivshort, Peng:24arxivshort}, GLA (a linear Transformer)~\citep{Yang:23arxiv}, HGRN and HGRN2 (RNNs)~\citep{Qin:23, Qin:24arxiv}, RetNet (linear Transformer)~\citep{Sun:23arxiv}, Hyena (linear Transformer)~\citep{Poli:23}, and additionally, xLSTM[1:0] and xLSTM[7:1], as defined in Section~\ref{sec:xlstm_blocks_notation}.

Training for these models employs mixed-precision techniques; however, in the cases of RWKV-5, RWKV-6, GLA, and HGRN2, PyTorch’s automated mixed precision was not used (refer to Appendix Section~\ref{sec:appGenTrainProc} for implementation specifics). The evaluated methodologies are organized into three principal groups: (a) Transformers, (b) State Space Models (SSMs), and (c) Recurrent Neural Networks (RNNs) in conjunction with linear Transformer models. Linear Transformers are distinguished as those leveraging linear mechanisms in lieu of the conventional attention component.

All architectures are scaled to match the GPT-3 parameterization at 350M parameters, featuring an embedding dimension of 1024 and 24 stacked residual blocks. Among these, only GPT-3 utilizes shared token and output embedding weights, resulting in a marginal reduction in parameter count relative to the others. Experimental results, presented in Table~\ref{tab:spaj15b_model_results}, establish that xLSTM achieves superior perplexity on the validation set compared to all other assessed models. Additional experimental and methodological details are provided in Appendix~\ref{sec:appExpComparison}. Figure~\ref{fig:temp_sclaw15B} illustrates the scaling trends, suggesting xLSTM’s advantageous performance persists as model size increases.

\paragraph{Ablation Studies.}
As shown in Table~\ref{tab:spaj15b_model_results} and Figure~\ref{fig:temp_sclaw15B}, xLSTM demonstrates strong performance in language modeling tasks when trained with 15B tokens from SlimPajama. In order to analyze the impact of each modification introduced in xLSTM compared to the standard LSTM, we incrementally transform a vanilla LSTM into an xLSTM architecture. The process starts by incorporating LSTM layers within pre-LayerNorm residual backbones. The next step involves broadening this design to include a post up-projection block. The final alteration consists of introducing exponential gating mechanisms along with matrix memory.

The results are shown in Table~\ref{tab:ablstudies} (top). 

The results from the ablation studies suggest that the notable gains in performance stem from the combined effects of exponential gating and the matrix memory mechanism. Furthermore, given the critical role of gating within RNNs and State Space Models, various gating strategies were systematically examined. As shown in Table~\ref{tab:ablstudies} (bottom), our analysis indicates that endowing each gate with input-dependent, trainable parameters incrementally improves model outcomes. Further investigation of the distinct backbone elements can be found in Appendix~\ref{sec:appAblDetails}.

\subsection{xLSTM as Large Language Model}
\label{sec:ExpLanguage}

This investigation concludes with extensive language modeling experiments designed to evaluate xLSTM as a potential large language model. To this end, we scale up the training set and utilize 300B tokens drawn from SlimPajama. Notably, this token count matches the setups adopted for models such as Mamba~\citep{Gu:24arxiv} and Griffin~\citep{De:24arxiv}.

We conduct a comparative analysis between xLSTM and three other models representing distinct methodological families discussed in Section~\ref{sec:ExpComparison}: RWKV-4, Llama, and Mamba. RWKV-4 is chosen as the RNN baseline because, for RWKV-5, RWKV-6, and HGRN2, appropriate precision configurations for training were determined only after launching the 300B token runs (refer to Appendix~\ref{sec:appGenTrainProc}). Our experimental protocol involves training the models at varying scales (125M, 350M, 760M, 1.3B parameters), evaluating their abilities to generalize to longer sequences, and measuring validation set performance. We further probe their effectiveness across downstream applications, assess language modeling proficiency on 571 text categories from the PALOMA benchmark, and analyze their scaling laws.

\paragraph{Sequence Length Extrapolation.}
\label{sec:spaj300B_ctx_extrapolation}
To begin, we evaluate how well 1.3B-parameter models—specifically, xLSTM, RWKV-4, Llama, and Mamba—generalize to longer input sequences. Although all models undergo training with a context length of 2048, their performance is assessed at test time for context lengths extending up to 16384. The outcomes of this analysis are depicted in Figure~\ref{fig:spaj300B_seq_extrapolation}. Notably, xLSTM demonstrates a superior ability to retain low perplexity as context lengths increase, outperforming the other approaches in this regard.

\paragraph{Validation Perplexity and Downstream Tasks.}
\label{sec:spaj_downstream_eval}
Next, across every model size, we assess how xLSTM, RWKV-4, Llama, and Mamba perform on the SlimPajama validation set with respect to next token prediction, as well as on downstream benchmarks designed to test common sense reasoning capabilities.

The third column in Table~\ref{tab:lmeval} presents the perplexity scores on the validation set for various approaches. Regarding validation set perplexity, both xLSTM[1:0] and xLSTM[7:1] consistently outperform alternatives across all model scales. Performance metrics for downstream tasks are given in the remaining columns of Table~\ref{tab:lmeval}. Across nearly all tasks and for every model size, xLSTM demonstrates superior results, with the exception of certain instances in the ARC task where Mamba achieves the best outcomes. Further information can be found in Appendix~\ref{sec:appExpLanguage}.

\paragraph{Performance on PALOMA Language Tasks.} Thirdly, across all examined model sizes, we evaluate xLSTM, RWKV-4, Llama, and Mamba models in terms of next token prediction accuracy on PALOMA language tasks~\citep{Magnusson:23arxivshort}. To quantify this, we use perplexity scores measured on 571 distinct textual domains, spanning sources such as nytimes.com and the r/depression subreddit. Table~\ref{tab:ppleval} presents perplexity values for next token prediction, separating results into broader language modeling tasks (first seven columns) and a selection of detailed domain-specific benchmarks (final five columns). On these language modeling assessments, xLSTM[1:0] exhibits superior performance compared to xLSTM[7:1].

xLSTM[1:0] achieves lower perplexity compared to Mamba across 568 of 571 (99.5\%) text domains, surpasses Llama in 486 of 571 (85.1\%) domains, and outperforms RWKV-4 in 570 of 571 (99.8\%) domains. Additional information is available in Appendix~\ref{sec:appExpLanguage}.

\paragraph{Scaling Laws.} Next, we investigate power-law scaling trends, which make it possible to project model performance as parameter counts increase~\citep{Kaplan:20,Brown:20short}. 
The scaling characteristics are depicted in Figure~\ref{fig:temp_sclaw300B}. Although the various models exhibit comparable scaling trends, they are separated by consistent offsets. Of the architectures compared, RWKV-4 shows the least favorable scaling, with Llama and Mamba following in succession. xLSTM outperforms Mamba, with the relative improvement roughly mirroring the difference between Mamba and Llama. These scaling patterns suggest that as model sizes grow, xLSTM is likely to maintain or even strengthen its performance lead over Transformers and State-Space models.

\textbf{Generation Times and Maximal Throughput.} In Figure~\ref{fig:inference_time_generation}, we present the generation times, and in Figure~\ref{fig:inference_time_decoding_throughput} (left), we display the peak throughput for various xLSTM configurations at the 1.3B parameter level. Our evaluation includes comparisons with equivalently sized implementations of Mamba, Llama, and RWKV sourced from HuggingFace. The Llama baseline also incorporates a static key-value cache. During these tests, neither full cache compilation of the Transformer Model nor the application of \texttt{torch.compile} to Mamba was successful. For text generation benchmarking, all models operate at a batch size of 1, with a pre-fill of 16 tokens—a setting that provides optimal conditions for the Transformer architecture.

The data in Figure~\ref{fig:inference_time_generation} highlights that xLSTM, alongside other recurrent approaches like Mamba and RWKV-4, exhibits linear complexity with respect to input length, in contrast to the quadratic behavior demonstrated by Llama. For decoding throughput, we conduct measurements across varying batch sizes and prefill settings, specifically examining the Llama model.

As depicted in Figure~\ref{fig:inference_time_decoding_throughput} (right), the constant memory requirements of xLSTM enable it to accommodate significantly larger batch sizes than Llama, ultimately resulting in superior throughput performance.

\section{Limitations}

(i) Unlike mLSTM, the design of sLSTM's memory mixing precludes parallel processing, thereby hindering rapid parallelization. Despite this, we engineered an efficient CUDA kernel for sLSTM, achieving a runtime less than double that of our parallel mLSTM implementation. (ii) It is important to note that the current mLSTM CUDA kernels have not undergone optimization, resulting in execution speeds that are approximately four times slower compared to FlashAttention or Mamba's scan operation. Performance gains similar to FlashAttention could be attained with more advanced CUDA kernel optimizations. (iii) Processing the $d \times d$ matrix memory in mLSTM introduces significant computational demands. Nevertheless, memory updates and reads are parameter-free and can exploit standard matrix-parallelization, keeping the actual time overhead due to the complex memory relatively low. (iv) Proper initialization of the forget gates is crucial to performance. (v) Given that the size of the matrix memory is sequence-length independent, increasing the sequence length risks saturating the memory with longer contexts; empirically, however, this has not constrained performance for contexts up to 16k tokens, as detailed in Section~\ref{sec:spaj300B_ctx_extrapolation}. (vi) Owing to the intensive computational requirements of large language model experiments, neither the architectures nor their hyperparameters were exhaustively tuned, particularly for larger xLSTM configurations. We expect that considerable architectural and hyperparameter optimization will be essential for xLSTM models to fully realize their performance capabilities.

\section{Conclusion}
Our investigation has provided a partial response to our primary question: What level of performance in language modeling can be achieved by scaling LSTM architectures to billions of parameters? Presently, our findings suggest: ``At least on par with state-of-the-art architectures such as Transformers and State Space Models''. We introduced xLSTM, a variant of LSTM incorporating exponential gating with memory mixing alongside a revised memory design. In language modeling tasks, xLSTM consistently demonstrates strong performance when compared to leading approaches like Transformers and State Space Models. Observed scaling laws suggest that increasing the size of xLSTM models positions them as strong contenders against today's prominent Large Language Models based on Transformer frameworks. Moreover, xLSTM presents promising opportunities for advancement in other domains, including Reinforcement Learning, Time Series Forecasting, and physical system modeling.

\end{document}